\chapter{Exactness and infinite discrimination}

\textbf{Every figure in this chapter is computed on the ideal filling and is deterministic; the runs
that produced them are named at each claim.} The ideal filling is Madelung order under Hund's rule,
stated in full in the next chapter. What is exact here is the arithmetic, and the figures measure what
it keeps.

\section{The engine reads by equality}

The anchor sift engine reaches a symbol through one operation. Its field interface answers, for two
positions in a corpus, whether they carry the same symbol: no order on the symbols, no distance
between them, and no alphabet they are drawn from. A measure built on that interface can read an
element, a letter or a chess move without knowing which, because all it ever asks is whether two
things are the same thing.

On integers that question is exact. Two integer quantum states are the same state or they are not,
and the reading never has to set a threshold for how close is close enough. That is the property this
book rests on, and it is worth stating plainly before any atom is read: where the symbols are exact,
the discrimination is unbounded. Nothing the source distinguished is lost, because nothing was
rounded into anything else on the way in.

\section{A coarse reading imposes a quantum}

The ingestion primitive keeps every digit its source wrote, and its header states the discipline it
holds to: no quantum is imposed from the reading end. A quantum is a decision about the smallest
difference worth telling apart, and the crystallography reader once made one silently, putting every
deposited site on a grid of a quarter of an angstrom, so that the recovered cell edge carried the
grid and not the deposit. The exact reading of the same cells answers the published edge with an
error of zero, and the difference between the two readings is the grid.

An atom has the same failure available. An electron sits at a quantum state of four numbers, the
principal, azimuthal, magnetic and spin numbers, and a reading that keeps only the first two, the
subshell, has dropped the two that tell the electrons of one subshell apart. The choice is the
voxel again, in a different domain.

\section{Pauli as a count}

The representation is the electron set of the neutral atom, one point per electron at an integer
state, built in
\texttt{examples/\allowbreak{}particle\_physics/\allowbreak{}1\_represent/\allowbreak{}an\_atom\_as\_exact\_shells.py}.
Two electrons at one state is a Pauli violation, and the same domain-blind primitive that finds two
elements on one crystallographic site finds it: \texttt{exact.contested} returns the positions
carrying more than one value. Over all 118 elements the exact reading places $Z$ electrons at $Z$
distinct states and no state is contested, so exclusion holds for every element the ledger names.

The check is proven by breaking its subject. A filling that placed two electrons at one state gives
\texttt{exact.placed}, a dictionary, one entry where the count expects two, so the distinct-state
count falls below $Z$ and the run reports it. A correct fill of two electrons keeps two states; a
doubled state keeps one. The count separates them, and the labels never disagree, because both
electrons carry the same subshell. The loss under a coarse reading is a count and only a count, and a
check that compares values reports nothing: across the table the subshell reading collapses 5645
electron states and a value-comparing check flags none of them.

\section{The partition is the quantum numbers}

The partition of an atom is which of the four numbers the reading keeps, and
\texttt{examples/\allowbreak{}particle\_physics/\allowbreak{}2\_partition/\allowbreak{}what\_a\_quantum\_number\_costs.py}
sweeps them. Keep all four and every electron stands at its own state. Drop the spin and the two
electrons of an orbital fall to one cell. Drop the magnetic number and a subshell falls together.
Drop the principal number and the atom is one cell.

The fullest one cell ever gets at each partition is the capacity Pauli allows there, and it is read
off the accumulation and not put in: one electron per spin orbital, two per orbital, fourteen at
the f subshell, thirty two at a full $n = 4$ shell. Only the full partition keeps every element's
electrons apart. Every coarser one agglomerates, and hydrogen, with one electron to fold, is the only
element that survives every partition. Exact arithmetic is the full partition, and the reading that
never agglomerates two states that differ is the reading this whole engine is built to be.
